\section{Balanced segments and continuity at all endpoints}\label{sec:endpoints}
The local comparisons of Section~\ref{sec:metric} identify the metric away from the
excluded points and give the deterministic scale estimates
\eqref{eq:original-18}. We now control endpoints uniformly. The proof first reduces
the problem to an estimate for bad times in balanced explorations. We then prove
the passage statement that enters the address count of Section~\ref{sec:address}.
All graph lengths below count every original primal edge, with its original unit
length; auxiliary Tutte paths serve only to locate the necessary primal joins.

\subsection{Balanced segments and the volume spine}

\begin{lemma}[A strict exponent margin]\label{end:exponent}
The critical metric exponent satisfies $\xi>1/4$.
\end{lemma}

\begin{proof}

We use the critical exponent characterization in the Ding--Gwynne LFPP results
\cite[Section~1.2, Equations~(1.4)--(1.7)]{DGuniqueness} (whole-plane pinned GFF, heat regularization at length \(\epsilon\),
median internal square crossing \(\epsilon^{1-2\xi+o(1)}\)). Indeed for any
\(\sqrt2<t<2\), all but \(o(\epsilon^{-1})\) of the \(\epsilon\)-squares in that square
have regularized field everywhere at least \(-t\log(1/\epsilon)\), in the sense that the
\textbf{number of exceptions} is \(o(\epsilon^{-1})\) in probability. This is by Gaussian
tails and first moment (use an intermediate exponent); unit-\(\epsilon\)-box field
oscillations have uniformly Gaussian supremum tails by the heat-kernel covariance
derivative bounds after scaling (center variance \(\log(1/\epsilon)+O(1)\)). Hence order
one of projected crossing length must use this lower height, giving \(1-2\xi\le t\xi\).

Since $t<2$, the inequality gives $\xi\ge 1/(2+t)>1/4$.

\end{proof}

\begin{definition}[Balanced segment]\label{end:balanced-segment}

In a Boltzmann disk $T_p$, start at a specified active edge without
choosing a target. At a face step choose the side fairly. At a bridge,
detach the smaller disk and retain the larger, with a fixed convention on
ties. Update the active edge to the nearest slot beyond it on the chosen
side, according to the two-list rule of Section~\ref{sec:cutting}.
Stop after
\[
 T=\max(1,\lfloor m_p\rfloor)
\]
steps, or on the first exit of the perimeter from $[p/2,2p]$.
For bounded positive $p$, choose $T$ bounded and take at least one step.
The segment children are all detached disks, including ring interiors,
and the final unexamined disk. Recurse in every child of positive perimeter;
a zero-perimeter singleton needs no segment.

\end{definition}

\begin{lemma}[Access to ancestral tip trajectories]\label{end:tip-access}
The balanced recursion exhausts the disk. Its active-edge midpoints can be
joined into tip trajectories $\gamma_p$ using boundedly many Tutte edge
portions per step. Every site can reach its ancestral tip trajectories by
following the recursion, with bounded local Tutte cost at each transition,
apart from travel along those trajectories.
\end{lemma}

\begin{proof}
Consecutive active edges have common or adjoining tip vertices: the nearest
retained slot shares the corresponding endpoint, and a bridge joins the
retained and detached endpoints. These incidences join successive active
midpoints, including the final absorption, at bounded Tutte cost.
For a ring child, start at the corresponding first inner base. It is
reachable by boundedly many Tutte links, since consecutive outer bases with
no intervening inner run adjoin the same inner vertex. Choose the root of
each recursive child at this transition site.

Every exposed site lies on a tip trajectory, on an adjoining ring or child
boundary, or in a singleton attachment. A site not yet exposed remains in
one of the children with its boundary occurrences retained. Iterating these
alternatives gives the asserted access to the ancestral trajectories.
The finite recursion exhausts the map by the exact peeling decomposition
\eqref{eq:original-5}.
\end{proof}

\begin{lemma}[Stopped word law]\label{end:segment-law}
For sufficiently large $p$, the complete word of a balanced segment in $T_p$ is dominated by a fixed constant times the raw iid word experiment of \eqref{eq:original-12}, with the same deterministic time bound and perimeter stopping rule. Conditional on the word, its unexamined fillings have the marked product laws.
\end{lemma}

\begin{proof}

Before stopping, the word of signed two-coordinate increments uses \eqref{eq:original-12},
with negative magnitude restricted to keeping the larger (\(j\le(S-2)/2\)) and the tie
factor, and overall likelihood \(f_{S_{\rm stop}}/f_p\); overshoot down is no smaller than
\((S-2)/2\), overshoot up unrestricted. In particular this whole stopped segment word for
large \(p\) is boundedly dominated by the raw iid experiment to time \(T\) with the same
stopping (invalid negative steps need not be used). Fillings are conditionally the marked
products as before.

The bound is uniform in the stopped word. Before exit, the perimeter is
comparable to $p$. A downward exit keeps at least $(S-2)/2$ and hence at least
a fixed positive multiple of $p$; an upward exit can be arbitrarily large, but
the perimeter weight decreases at the polynomial rate already established.
The ratio $f_{S_{\rm stop}}/f_p$ is therefore bounded in both cases. The
restriction on negative jumps and the tie convention can only remove raw words
or decrease their multiplicity.

\end{proof}

\begin{lemma}[Volume weights and logarithmic spine moments]\label{end:volume-spine}
Let $V(p)=\E N_p$. Then
\[
 V(p)\sim\frac{c_*^2}{4f_p}\asymp\frac{p^2}{\log p},\qquad p\longrightarrow\infty,
\]
with $V(p)>0$ for $p>0$ and $V(0)=0$. For the segment-child kernel
\[
 \mathcal K(p,q)=\frac{V(q)}{V(p)}
 \E\#\{\text{children of perimeter }q\},\qquad p,q>0,
\]
volume additivity makes $\mathcal K$ substochastic. As $p\to\infty$, its mass tends
to one, its logarithmic ratio drift tends to zero, its logarithmic ratio second
moment stays bounded above and below by positive constants, and its logarithmic
ratios have uniformly bounded two-sided exponential moments of sufficiently
small order.
\end{lemma}

\begin{proof}

In \eqref{eq:original-6} the perimeter moment's \(c\)-derivative equals \[
 p F_{p-1}+\tfrac14\,\mathbb E[(c-U)^p(U-\eta/U)]
\] by the differentiation above. At \(c\downarrow c_*\) the \(\eta/U\) term is bias
converging to an atom of mass \(c_*\) at zero (\(\eta=c/m_{-1}\)). Consequently the
derivative divided by \(c_*^p\), after this limit then \(p\to\infty\), tends to \(-c_*/4\)
(the other terms vanish by the moment equivalent and the endpoint gap). Thus
\(V(p)\sim c_*^2/(4 f_p)\asymp p^2/\log p\), finite positive for \(p>0\), zero for
\(p=0\).

Use the substochastic spine kernel on positive perimeters \[
             {\cal K}(p,q)=\frac{V(q)}{V(p)}
                    \mathbb E\,\#\{\text{segment children of size }q\},
\] using volume additivity/domination. Its large-\(p\) masses tend to one, log ratio drift
to zero, variance stays between positive constants, with uniformly bounded two-sided small
exponential moments of log ratio. Here is the stopped-segment calculation. Let $\tau$ be the stopping
step, $S_r$ the retained perimeter just before step $r$, and $Q_r$ the
perimeter of the child detached at that step (the inner perimeter at a
ring step). The final retained disk is itself a child. For $z$ near $2$,
with $0^z=0$, the exact telescoping identity is
\[
 \E\sum_{q\text{ child}}(q/p)^z-1
 =\E\sum_{r<\tau}p^{-z}
       \bigl(S_{r+1}^{\,z}+Q_r^{\,z}-S_r^{\,z}\bigr).
\]
The first and second derivatives in $z$ give the corresponding
logarithmic-ratio moments. We estimate the one-step summands before
summing them over the stopped path. For one step from
\(s\to\infty\), the scaled drift (including the detached child) of perimeter power \(z\)
relative to \(s^z\), times \(m_s\), tends near \(z=2\), with first two derivatives, to \[
       G(z)=\mathrm{PV}\int_{-1/2}^{\infty}
       \frac{(1+v)^z-1+|v|^z}{(1+v)^2}\frac{dv}{v^2}.
\] Indeed use the \eqref{eq:original-12} equivalents, \(f_{s+d}/f_s\) likelihood, and the
positive-face inner/outer ratio tending to one with geometric-run and smooth tilt tail
bounds. Truncated signed first moment of the unrestricted raw increments is bounded; at
small \(|d|/s\) use \(|f_{s+d}/f_s-1|\le C|d|/s\), second-order cancellation for the
survivor term and absolute power integrability for the detached term. For the latter all
positive outer sizes \(k\) have inner powers up to slightly past 2 bounded in expectation
by \(C(1+k)^z\) from the geometric sum tilt; use nearby powers also for log factors. At
the upper tail the additional likelihood decays as at least power \(2-\delta\); these
bounds give dominated truncations. 

For the small-step likelihood bound used here one may
take the limiting positive moment law from \eqref{eq:original-6}, normalized in \(c_*\)
units, whose negative endpoint has modulus \(<1\). On \(y\in[0,1]\),
\(|y^{s+d}-y^s|\le C (|d|/s)y^{s/4}\) for \(|d|\le s/2\); moments there at comparable
indices are bounded by \(O(f_s)\) by the equivalent and the exponentially small
negative-cut contribution. The latter itself causes no problem for nonzero integer \(d\).
The two coordinates together have the stated unit Lévy coefficient. 

Now \(G(2)=0\) by
integrating \(2/[v(1+v)]\). For \(G'(2)\) the first integral
\(\mathrm{PV}\int \log(1+v)dv/v^2=2\log2\); the second is
\(\int \log|v|\,dv/(1+v)^2=-2\log2\) by substituting \(v/(1+v)\). Also \(G''(2)>0\). 

Sum
these drift estimates over pre-stopping states (all comparable to \(p\)), multiplying by
\((s/p)^z\) and including derivative terms; expected step count divided by \(m_p\) is
bounded and bounded below (a small initial time with small raw oscillation has positive
probability by Cauchy convergence and the bounded positive likelihood there). This gives
the claimed weighted estimates first with \((q/p)^2\) weights. Replacing by \(V(q)/V(p)\)
is harmless by the nearby power bounds and the equivalent. Indeed the ratio is comparable
on compact positive ratio ranges and controlled off them by arbitrarily small
positive/negative power losses relative to \((q/p)^2\) from the
\(\asymp (1+p)^2/\log(2+p)\) estimates on positive perimeters, with relative multiplier
tending to one on the compact ranges.

\end{proof}

\begin{proposition}[Simultaneous segment counts]\label{end:band-counts}
Start a balanced recursion at $p_{\rm init}\asymp B_n$ and stop a branch when it
first exceeds $p_*=C_*B_n$. The probability of any such crossing tends to zero
as $C_*\to\infty$. For every $\delta_1>0$, sufficiently small, the expected
number of nodes with $k=\lfloor\log(p_*/p)\rfloor$ is at most
$C e^{(2+\delta_1)k}$. With arbitrarily high probability, simultaneously for all
ancestry lines and all $k\ge0$, the number of such nodes on any one line is at
most $C e^{\delta_1k}$, with a suitable deterministic $C$.
\end{proposition}

\begin{proof}

Start at \(p_{\rm init}\asymp B_n\) and truncate on rising above \(p_*=C_* B_n\), which
somewhere in the tree has probability bounded by \(V(p_{\rm init})/\inf_{p>p_*}V(p)\) by
the spine/first-hit sum; \(C_*\) can be chosen large. Put
\(k=\lfloor\log(p_*/p)\rfloor\ge0\) for bands. For each small \(\delta_1>0\) the expected
number of nodes in band \(k\) before such crossings is \(\le C\exp((2+\delta_1)k)\).
Moreover with arbitrarily high probability, simultaneously over all ancestry lines and
\(k\), at most \(C\exp(\delta_1 k)\) such nodes occur per line. To see it on the
killed/upper-stopped spine, tilt steps at large \(p\) by \(e^{\theta\log(q/p)}\) and
normalize for fixed sufficiently small \(\theta>0\). Normalizers then exceed one, tilted
drift is bounded below positively, and a small negative exponential of the increment has
tilted expectation \(<1\), uniformly beyond a fixed cutoff, by the preceding moment
bounds. From a band visit sufficiently high above that cutoff, with uniformly positive
tilted probability the chain stays above cutoff and crosses the upper barrier in
\(O_\theta(k+1)\) steps with log total rise \(\le C_0(k+1)\) at crossing (absolute
coefficient for large \(k\)). Indeed use the negative-exponential supermartingale for
avoiding cutoff, the positive drift and bounded variances, and the step tail for the
overshoot. Thus before tilting a kill/crossing within that many steps still has
probability at least \(c_\theta e^{-C_0\theta k}\). On the remaining fixed bounded
perimeters there is uniform direct kill probability (positive expected exposed ring
volume). Let $N_k$ be the number of visits to band $k$ by this killed,
upper-stopped spine. Restarting at later visits gives
\[
 \Prob_{\rm spine}(N_k\ge m)
 \le \exp\!\left[-q_k
          \left\lfloor\frac{m}{C_\theta(k+1)}\right\rfloor\right],
 \qquad q_k\ge c_\theta e^{-C_0\theta k}.
\]
The volume bias converts this to both asserted tree estimates. Indeed,
\[
 \E\sum_{v:\,p_v\text{ in band }k}
          \frac{V(p_v)}{V(p_{\rm init})}
       =\E_{\rm spine}N_k.
\]
On this band, $V(p_{\rm init})/V(p_v)\le Ce^{2k}$.
The same identity at the stopping line of the $m$th band visit gives
\[
 \Prob\{\text{some ancestry line has at least }m\text{ band visits}\}
 \le Ce^{2k}\Prob_{\rm spine}(N_k\ge m).
\]
Choose $\theta$ so small that $C_0\theta<\delta_1$ with strict slack.
Summing the geometric tail proves the expected-node bound. Substituting
$m=Ce^{\delta_1 k}$ and then summing over $k$ proves the simultaneous
ancestry bound, with arbitrarily high probability as $C$ increases.
Fresh
disk tests at a node can then use its unconditional \(T_p\) probabilities.

\end{proof}

\begin{lemma}[Transfer and local incidence costs]\label{end:finite-sphere-transfer}
The preceding recursion estimates apply, with bounded domination, to both
initial disks in a restricted central comparison with the unit finite sphere.
Under the rooted sphere law, the maximum vertex degree and maximum face degree
are $O(\log n)$ with probability tending to one. On this event a bounded local
Tutte advance or incidence join can be replaced by a primal join of cost
$O(\log n)$.
\end{lemma}

\begin{proof}

Both initial disks of the restricted central comparison obey these bounds with bounded
domination in the unit finite sphere: leave the opposite disk for volume smoothing by
\eqref{eq:original-10}, localizing both initial perimeters in positive compact ranges.
Also the full sphere has vertex and face degrees \(O(\log n)\) with probability tending to
one under the rooted law. Indeed at a vertex of root degree \(d\) split its list from the
root into \(j+1\) nonempty consecutive groups in order, replacing the vertex by a path
with \(j\) occupied connecting edges and placing the groups along one flank at the
successive vertices with the other flank free. Given \(j\) invert by tracing that empty
flank starting adjacent to the root half-edge in the specified rotation and contracting;
the FK exponent is unchanged. Thus a root-origin degree tail costs at most
\(Z_{n+j}/(Z_n\binom{d-1}{j})\); use \eqref{eq:original-3} and \(j\) a sufficiently small
fixed fraction of the threshold for exponential decay. Use duality and uniform rerooting.
This high-probability bound may be intersected \textbf{after} transferring estimates back
from central intensity by dividing by counts.

On the degree event, incident primal vertices around a face have primal
joining distance $O(\log n)$. Within a disk we use genuine primal edge portions,
including the possible blue half-diagonals up to a boundary red base; the costs
of these internal primal walks carry to the joined sphere.

\end{proof}

\subsection{The bad-time estimate and the endpoint reduction}

We give the reduction to a one-segment estimate, stated explicitly to separate the
endpoint issue. Choose \[
 1/4<\tau<\varrho<s<\sigma<\min(\xi,1).
\] Use \eqref{eq:original-18} with \(\sigma\). Take \(b>1\) sufficiently close to 1
(\(s-(b-1)>\varrho\)), then \(a>1\) much closer to 1 with \(a<b\). A small truncation
constant \(\delta>0\) can first be chosen so that \(\sigma-\varrho\gg\delta\), then
\(b-1\ll\delta\). Consider diverging perimeter parameters \(p_*,p\), put
\(k=\lfloor\log(p_*/p)\rfloor\), and take integer \(N\to\infty\) with \(k\le C N,\ k+N\le(1-\delta)\log p_*\). Time sizes will be
\(T e^{-x}\). In each time interval of size \(T e^{-N}\) test grid times at spacing
\(T e^{-bN}\) (integer rounding permitted). Only test times of the executed path outside
endpoint strips of width \(O(T e^{-aN})\) at start and stop.

\begin{definition}[Good grid time]\label{end:good-time}

A good time has a simple short separating boundary, either a primal circuit in the
disk or a simple primal crosspath capped in the external root face. Its real portion costs
at most \(\mathcal B(p_*) e^{-s(k+N)}\) up to harmless bounded factors. One open side
contains the path throughout a time neighborhood of radius at least the grid spacing
(rounding with slack); both full-segment ends are strictly outside, and the component time
interval in that side about the tested time ends on both sides within \(T e^{-aN}\) up to
fixed factors. Shifted constants/finer intermediate ranges with slack are harmless.

\end{definition}

\begin{proposition}[Bad grid times]\label{end:bad-grid}

For each prescribed \(B,C\), with failure probability \(O(e^{-B N})\) under \(T_p\),
\textbf{every} such base interval (intersect the executed path) has at most
\(\exp(\tau(b-1)N)\) bad grid times, outside the stated strips. It suffices to use a mesh
of base intervals and allow a constant multiple of this bound on arbitrary intervals. In
this reduction and the estimate one can spend strict exponent slack in all the displayed
inequalities.

\end{proposition}

Proposition~\ref{end:bad-grid} is proved in Section~\ref{sec:address}. The
remainder of this section prepares its passage input. Before doing so, we show
exactly how that proposition completes the missing endpoint control.

\begin{proposition}[Uniform endpoint continuity]\label{end:endpoint-reduction}
Assume Proposition~\ref{end:bad-grid}. For every $\varepsilon>0$,
\[
 \lim_{r\downarrow0}\limsup_{n\to\infty}
 \Prob\left[
 \sup_{\substack{u,v\in V(M_n)\\\rho(\phi_n(u),\phi_n(v))\le r}}
 \frac{d_n(u,v)}{\mathcal B(B_n)}>\varepsilon
 \right]=0.
\]
Normalized diameters are tight. Every subsequential limit of the entire embedded
distance graph equals the same deterministic multiple of $D_h$ already
identified by the local open-port comparisons, including at the excluded points.
\end{proposition}

\begin{proof}

\smallskip
\noindent\textbf{Intersecting good intervals and the recurrence.}\quad

Indeed overlapping good component intervals with neither contained in the other have
intersecting real boundary paths. For two capped paths, if cap endpoints alternate, the
real paths already have to meet by planar order (the root face gives the occurrence order
even when the boundary is pinched); otherwise one can cap disjointly when testing that
pair. Coincident real endpoints already give contact. Nonintersecting boundaries would
then be nested if their open chosen sides overlap, since both avoid the common start point
also in their closures, and the overlapping component intervals would nest. Internal
capped/full contacts suffice anywhere on the real portions. Thus one can prune
containments in each intersecting chain and join along actual paths paying at most their
total real lengths and the local incidence corrections. The endpoints at uncovered gaps
are actual hits on those paths. Uncovered gaps across a tested interval require in sum
\(O(1+\#\mathrm{bad})\) intervals of size \(O(T e^{-bN})\), since every non-end good grid
time covers a spacing on both sides. More explicitly gaps of even shorter length must have
separating bad centers between the good centers of their two chains. Endpoint remainders
caused by using intersecting intervals, or by skipping the two full-segment end strips,
cost only \(O(1)\) intervals of size \(O(T e^{-aN})\).

Consequently if \(\Omega(N)\) bounds joining costs for tip-point pairs within \(T e^{-N}\)
(using adjoining blue representatives in the joined sphere), it suffices recursively to
pay
\[
 \begin{aligned}
 &O(1)\Omega(\lfloor aN\rfloor)
   +O(e^{\tau(b-1)N})\Omega(\lfloor bN\rfloor)\\
 &\qquad+O(e^{(b-1)N})
       \bigl[\mathcal B(p_*)e^{-s(k+N)}+O(\log n)\bigr].
 \end{aligned}
\]
The second line accounts for circuits, capped paths, and incidence
transitions. Fixed size factors can instead shift log indices by bounded amounts. Start
this on all integer levels \(N\ge\max(M_0,c_2 k)\) with small \(c_2>0\), stopping
recursion at \(k+N\ge(1-\delta)\log p_*\). In the terminal range reached from below
(\(k+N\le(1-\delta/2)\log p_*\) after choices), trivial chronological costs
\(O((1+T e^{-N})\log n)\) are negligible relative to \(\mathcal B(p_*)e^{-\varrho(k+N)}\).
For completeness, divide each recursive term by the proposed bound
\(Q(N)=\mathcal B(p_*)e^{-\varrho(k+N)}\). Apart from bounded shifts of the indices,
the two recursive coefficients are bounded by
\[
 C e^{-\varrho(a-1)N}
 \quad\text{and}\quad
 C e^{-(\varrho-\tau)(b-1)N},
\]
and the circuit-cost ratio is at most
\[
 C e^{-(s-\varrho)k-[s-\varrho-(b-1)]N}.
\]
Each tends to zero by the strict exponent choices. The incidence errors are
absorbed by \(b-1\ll\delta\ll\sigma-\varrho\). At the terminal levels the
chronological time is at most a constant times \(p_*^\delta\), up to logarithmic
factors, whereas \eqref{eq:original-18} makes the proposed bound a larger
positive power of \(p_*\); the same margin absorbs the terminal \(\log n\) costs.
Choose \(M_0\) so that the sum of the recursive coefficients is less than one,
and use the negligible terminal costs to initialize downward induction. This
gives the bound \(Q(N)\) throughout the recursion. Bands starting
already too small need only trivial whole-trunk costs, \(\ll\mathcal B(p_*)e^{-\eta_0 k}\)
for some \(\eta_0>0\) by taking \(c_2,\delta\) small. Summing starting-mesh intervals on a
trunk similarly yields \(\ll_{M_0}\mathcal B(p_*)e^{-\eta_0 k}\), as well as normalized
chronological equicontinuity on bounded \(k\). By the spine bounds and fresh-disk
estimates with \(B\) large, \(M_0\) can give arbitrarily high joint success probability
over nodes and levels. Sum along each ancestry using \(\delta_1\ll\eta_0\), also absorbing
\(O(\log n)\) attachment errors per node. This implies uniform approximation of vertices,
at normalized cost error tending to zero with band cutoff, by points on the finitely many
macroscopic tip trajectories (use the last sufficiently large ancestor). Central ring
sites are covered on the initial disks. In these high-probability arguments
\(\mathcal B(p_*)/\mathcal B(B_n)\) is bounded for fixed truncation by the adjacent-ratio
bounds. One may first spend an arbitrarily small \textbf{fixed} probability error on the
perimeter/degree truncations and the constants in the visit and trunk bounds, then test a
putative failure in a further coupled limit on the remaining events (with bounded-band
node counts tight by the spine estimate, or truncated simultaneously by growing
band-dependent cutoffs). Thus no almost sure constant bound for untruncated maxima through
all \(n\) is needed.

\proofstage{Nonconstant limits of macroscopic tip trajectories}

The preceding estimates approximate every vertex, with arbitrarily small
normalized cost, by a point of one of finitely many macroscopic tip trajectories.
We need one more property of those trajectories: a limiting trajectory cannot
remain at one location throughout a nonzero time interval.

First extract the macroscopic trajectories jointly, parametrized by their
rescaled chronological times. Cost equicontinuity and the positive lower
comparison on ordinary compact advances give spatial equicontinuity. Indeed,
a fixed spherical advance contains a compact subadvance avoiding the finitely
many excluded points, and that subadvance has a positive normalized cost.
The duration of a macroscopic segment does not collapse in probability:
its stopped word has bounded density relative to the raw word, and the
Cauchy limit stays in the stopping range for a positive initial time.

On every nonzero time interval before stopping, positive jumps of the raw
Cauchy coordinates are dense. Choose two distinct positive jumps in that
interval. Their limiting sizes are positive, and the ring comparison aligns
their inner and outer perimeters. By the disk-volume limit
\eqref{eq:original-9} and conditional freshness, both born disks have positive
limiting area. Each also leaves positive exterior area, supplied, for example,
by the other disk. The area comparison and diffuseness then prevent either
corresponding FK loop from shrinking to a point. The tip is incident to each
loop at its birth. Distinct macroscopic limiting CLE loops are disjoint,
so the tip limit cannot be constant on this interval.

These statements are first made for finitely many strict surviving tests.
Exhaustion over such tests, the bounded-band spine counts, and domination by
the unobserved opposite initial disk make them simultaneous for the finitely
many trajectories retained in any fixed truncation. Taking rational time
intervals proves nonconstancy on every nonzero subinterval of each retained
limit trajectory.

\proofstage{Short surrounding circuits and the endpoint contradiction}

There are surrounding primal circuits of arbitrarily small normalized cost
with strict spatial clearance about every sphere point. For ordinary points,
use the aligned estimate with reference-metric tests: the length-space topology
supplies finite-length winding and transverse crossing paths in a prototype
annular buffer, and the unit actual-height comparison supplies shrinking
successful annuli. To include the two marked ends, resample the area-marked
representation. Almost surely the old ends are ordinary points in the new
representation. At each fixed test the reference paths can avoid all excluded
points, so the open-port upper comparison produces the required primal
circuits. The same construction gives short reference-metric bypasses of
the excluded points.

We prove the asserted endpoint modulus by contradiction, making the order
of the truncations explicit. Suppose there are $\varepsilon_0,\eta>0$, a
sequence of map sizes, and vertex pairs $(u_n,v_n)$ such that
\[
 \rho(\phi_n(u_n),\phi_n(v_n))\longrightarrow0,
 \qquad
 \Prob\{d_n(u_n,v_n)>\varepsilon_0\mathcal B(B_n)\}>\eta.
\]
Measurable choices are available since each vertex set is finite. Fix the
perimeter and degree truncations and the spine/trunk constants so that their
total failure probability is less than $\eta/2$. Next choose a band cutoff
so that, on these events, each vertex can reach one of the retained
macroscopic trajectories at normalized cost less than $\varepsilon_0/8$.
Include the failure indicators, the chosen trajectory indices and joining
times, and the connection costs in this extraction. Pass to a further
joint coupling of these variables, the finite trajectory family, and the
chosen endpoint locations. The
intersection of failure with the retained truncation events has limiting
probability at least $\eta/2$. On that limiting event the endpoint images
approach the same point $z$, and the failure indicators are eventually one.

Consider the low-cost connection from $u_n$ to its approximating trajectory.
If its terminal point has limiting location different from $z$, this
connection already reaches away from $z$. If its terminal point limits to
$z$, move a sufficiently small fixed chronological distance along that
trajectory. Nonconstancy supplies a different limiting location, and cost
equicontinuity makes the added normalized cost less than
$\varepsilon_0/8$. Apply the same argument to $v_n$. We obtain two primal
connections, each of normalized cost less than $\varepsilon_0/4$, from the
approaching endpoints to points whose limiting locations differ from $z$.

Choose a sufficiently small surrounding circuit about $z$ which excludes
both of these latter locations and has normalized cost less than
$\varepsilon_0/4$. Its strict clearances force both connections to meet its
real primal path for all sufficiently large $n$. Joining at those
intersections gives
\[
 d_n(u_n,v_n)<\tfrac34\varepsilon_0\mathcal B(B_n),
\]
contradicting the failure. All choices were made after fixed truncations
losing less than the putative failure probability. Letting that truncation
error decrease proves the endpoint modulus in probability. This argument
uses no almost-sure bound on untruncated maxima through all map sizes.

\proofstage{Diameter tightness and identification of the whole distance graph}

The endpoint modulus, finite spatial chaining, and vanishing flag mesh give
tightness of normalized diameters. The open-port upper comparison follows
finite-length reference paths away from the excluded points. Endpoint
continuity attaches the specified vertices, while the short bypasses and
the reference-metric completion remove the excluded points. This gives the
upper bound by the same deterministic multiple of $D_h$ for every pair.

For the lower bound, bounded-cost graph geodesics have continuous limiting
projections after extraction: every fixed spatial advance has a positive
lower charge. Off the excluded points their compact subpaths obey the local
lower comparison. Portions visiting a shrinking neighborhood of an excluded
point can be omitted at vanishing reference-metric error. More explicitly,
skip from the first to the last visit to a neighborhood of the first such
point encountered, and then do the same for the other point if it is
encountered later. The triangle inequality and continuity of the completed
reference metric bound the resulting reference error by the diameters of
these neighborhoods. Let those diameters tend to zero. The lower and upper
bounds therefore agree on the entire embedded distance graph.

\end{proof}

\subsection{Tangent tests at a balanced active slot}

We prepare the proof of Proposition~\ref{end:bad-grid}. Time indices refer
to states (between moves); bounded roundings, and parametrizing a move including its
endpoint transitions within a bounded extension into adjacent step intervals, will cause
no problems. Throughout put \[
 r_y=T e^{-y},\qquad P_y=(2/\pi^2)r_y\log r_y.
\] These units are only used for \(0\le y\le bN+O(1)\). Thus \(m_{P_y}\sim r_y\), the
logarithm in \(P_y\) is comparable throughout to \(\log p_*\),
\(P_y\asymp_\delta p e^{-y}\), and the smallest \(r_y\) here is greater than a fixed
positive power of \(p_*\). Bounded shifts of \(y\) have their ordinary exponential scale
ratios asymptotically. All constants/buffers can depend on fixed exponents and (when used)
on a fixed moment order, not on the large parameters. Fix also a little shorter scale
range \[
                 {\cal I}=[a_+ N,b_- N],\qquad 1<a<a_+<b_-<b
\] with \(a_+-1\) and \(b-b_-\) sufficiently small compared with \(b-1\) and our strict
exponent slacks. On this range a budget of any fixed multiple of \(\mathcal B(P_y)\) is
negligible compared with the good-time budget by \eqref{eq:original-18}.

\begin{lemma}[Literal persistence of ordinary slots]\label{end:slot-persistence}
A lifted open interval of pending ordinary labels, strictly separated from all
thresholds swept by intervening operations, persists with its exact slot
identities, boundary order and color. A strict interior interval in an inserted
positive aperture is new ring data; one in a removed negative aperture passes
literally to the detached disk. The assertion also applies to transfer into a
survivor, and it preserves the prescribed corner convention through a ring.
\end{lemma}

\begin{proof}

Describe the main balanced segment's pending ordinary list (all its boundary slots except
active) by one lifted interval with successive integer vertices \([b_u,h_u]\), of width
equal to perimeter minus 1; the active slot reconnects its two ends. We use one oriented
order consistently. The walk is in \((h_u,-b_u)\); exactly one coordinate updates by
\eqref{eq:original-12}. On a positive expansion step (\(k-1\) replacement slots), \(k-2\)
is added on the chosen side, leaving the last slot at the other end active; \(k=1\) is the
previously described deletion/reassignment. On a negative at one end the \(j\) intervening
slots there pass into \(T_j\), and \(j+2\) ordinary slots disappear at that end in the
main list (reverse-glue, reassign). Thus an open interval of labels strictly between both
end thresholds before and after moves that do not sweep through it is unchanged, literally
slot by slot. On a positive step it can arrive entirely new from a ring if strictly inside
the inserted range; on a negative step it passes literally to the detached list if
strictly inside the removed range. Here and below strict clearances in assertions at scale
\(P_y\) absorb the bounded endpoint conventions. Along increasing labels the order
relative to the ring, or to a detached list, is the same whichever end made this update:
use the direction around the oriented unfilled polygon on ordinary slots. The ring inner
boundary has its fixed gluing order. Color across a positive changes and into a negative
does not. These conventions also apply to arcs of the initial boundary and in a survivor
at segment stop. In particular gluing at consecutive ordinary slots transferred into a
bridge child or survivor uses their inherited boundary order with no new phase on that
arc. New ordinary outer bases at either end of a face update can be read in the same
increasing frontier order along an open inserted interval, paired through binary runs to
the oppositely glued inner disk. Chunks strictly inside such an interval can always take a
cut at a specified one of these outer bases and the same incident inner-corner convention,
irrespective of which end of the main list supplied the insertion. The iid geometric gap
calculation before tilt works in this cyclic convention as well.

\end{proof}

Use a spatial tip path for instance walking between the midpoints of successive actives
via the shared endpoint on reassignment/insertion as above, recording the other
peeled/attachment endpoints if desired by bounded returns. Its movement uses adjoining
Tutte edges. It is internal to the disk including its boundary walk. At all its visits it
is within bounded local incidence of the active edge whose pairing is made or new face
exposed in that operation (use the last operation at stop). This variant works in the
ancestry argument: transition sites at a split or a ring pinch can use boundedly many
adjoining incidences, and actual boundary vertices not on the swept tips continue into
children. Root edges for the recursive balanced segments of children are taken at those
transition pinches/tips (or a nearest inner base with the bounded Tutte access). This also
covers initial-boundary sites in the sphere assembly.

\subsubsection{The lifted bad-address law and helper prescriptions}

Suppose contrary to the estimate along a large-parameter sequence that one interval of
time width \(O(Te^{-N})\) has \(>\exp(\tau(b-1)N)\) bad grid times with probability at
least \(\exp(-B' N)\). Here and below only actual times being tested are counted. It
suffices to disprove this for arbitrarily large fixed \(B'\), absorbing the number of
intervals. Fix a large integer \(d\) (moment order). On a new law condition on this event
and sample \(d\) such bad times with replacement in the interval; call these the bases
\(t_i\). They lie strictly on the executed path with both endpoint distances large
relative to all bounded windows at \(y\in{\cal I}\). Extend this change of law to all
auxiliary observations by using their given conditional laws. Relative to disk law with
\(d\) uniform time-grid choices in that interval it has relative entropy \(O_{d,B'}(N)\).
Relative to disk law times \emph{counting measure} on the \(d\) addresses it has density
bounded by \begin{equation}
\exp\{ B'N-d\tau(b-1)N\}.             \label{eq:original-19}
\end{equation} We refer to the new law as the lift. Extra harmless randomness, such as
grid phases, may be added independently. Our reductions are on further subsequences. The
entropy constant need only be finite.

Here is some auxiliary exploration we can use in this test. First generate the whole main
word and its aperture/ring parameters. In its individual offspring disks \(J\) (of either
sign, also the final survivor if any), preassign finitely many boundary labels, at most
two per base: use \(h_{t_i},b_{t_i}\) if in the corresponding lifted interval, transported
through the ring order if positive. Interior positions may be rooted at an adjoining slot
by one fixed convention, including nearest inner base by the ring list. Endpoint or
degenerate cases can use arbitrary clipping or omission when the label is not in ordinary
transfer. These seeds do not inspect the filling. An original boundary occurrence is \emph{consumed} when it is reached as
active or used as a paired or deleted slot. Passing an untouched interval
into a fresh detached child is transmission, not consumption. In each
\(J\), treat the prescribed seeds successively in a fixed order. Use one balanced segment from the prescribed slot if still pending in a
fresh polygon, leaving all its child and survivor fillings unexamined for further
prescriptions; skip already consumed/unavailable slots. The prescriptions concern original
boundary occurrences, carried down the splits or main changes until consumed. Earlier
operations here don't fill polygons unnecessarily. Thus on the untilted law each started
trajectory, given its starting polygon size and the preceding word histories, has the same
bounded iid dominance (extend by raw increments after stopping), and fresh raw detached
contents when those alone are queried along it. One need not treat its trajectory as
independent of an observed interior of that same disk. Every untouched arc of consecutive
original slots with a slot margin around it lies in a single such polygon until affected:
a change separating or replacing part of that arc has to reach an active or paired slot
there. Faces emit at, not between, the active slots. Cyclic order and the exact slot
identities persist under this transmission.

\subsubsection{Typical-scale tangent laws}

Use a large integer shell ratio \(R\), \(g=\log R\), fixed first. Sample \(y\) uniformly
on \(g\mathbb Z\cap{\cal I}\). The following local blow-up observations under the lift
will be used.

\begin{lemma}[Typical-scale absolute continuity and core alternatives]\label{end:typical-tangents}
\leavevmode

\begin{itemize}
\item
  Around a base \(t=t_i\), increments on compact time subintervals of
  \(\mathbb R\setminus\{0\}\), after translating time \(t\) to zero and rescaling by
  \(r_y,P_y\), have subsequential limits absolutely continuous there (for increments,
  relative to the two raw Cauchy coordinates). Include also maximum oscillations from the
  center in symmetric windows, in the extended half line. On the event these are finite in
  the limit, joint relative positions measured from the center off time zero are
  absolutely continuous on compact value ranges relative to smooth initial-position
  densities times the raw increment laws. Use integer-time paths with side copies as in
  the \(J_1\) comparison. The limit can be taken stationary under shifting \(y\) by \(g\),
  jointly at all fixed such shifts.
\item
  The centered path thus either has infinite core oscillation in every positive-size
  window (a singular test) or those oscillations tend to zero as window size tends to
  zero. Indeed off-center increments are finite, so the infinite case cannot start at a
  positive radius; in the finite case use monotonicity and dilation stationarity. On the
  finite event, for each of the two coordinates in each of the two directions from the
  center, either it goes strictly below the center height somewhere at finite distance, or
  it stays strictly above it on all compact time ranges in that open half-axis (infima
  include left limits). Nonnegative touching with a zero infimum on a rational inset
  interval is excluded by initial-position absolute continuity. We call the latter cases
  positive halves.
\end{itemize}

\end{lemma}

\begin{proof}

To see the entropy assertion precisely, space out annular blocks \([\eta,U]\) and
\([-U,-\eta]\) in time units across \(y\) in finitely many sublattices so they are
disjoint for fixed \(0<\eta<U<\infty\), leaving also intervals of length comparable to
\(r_y\) between consecutive blocks on both sides. Fix the proposed bases initially. The
full augmented main word relative to raw iid (continued if necessary) has the bounded
likelihood from the balanced rule. Interior constraints on observations can simply drop a
query when outside the word's executed range, or include artificial raw outputs. Read
these raw word blocks from inner to outer, including their positions at the near endpoints
relative to the base center (using a cutoff symbol when outside a desired compact range).


Sums in the unused fresh intervening gaps have local mass bounded by \(C/P_y^2\) each for
the two coordinates. Indeed condition on the binomial fair-side counts, ignoring
exponentially unlikely imbalance; the two sums are then independent with the scalar
lattice stable local limit bound (span one, the tail equivalents and bounded truncated
mean of \eqref{eq:original-12}). Positions for the two directions use independent gaps.


Thus these outputs are sequentially dominated, with bounded cost in relative entropy per
queried scale, relative to raw block words with independent uniform lattice positions on
the prescribed bounded boxes (plus a cutoff symbol of positive mass). The chain rule,
mixing over scales and the \(O(N)\) budget give bounded relative entropy at a typical
scale. In particular any tests null asymptotically under those comparison references are
null here. This gives the asserted absolute continuity and tightness using the walk
convergence and increasing cutoffs, including on finite core-oscillation events. The
lattice spacing here is negligible; independence of initial positions in the limit
assertion is required only on the reference side. For a path fixed not to extend to both
directions one can first adjoin raw data past its ends. 

Stationarity follows by the
uniform log-level shift (and exact asymptotic dilation factors) in diagonal extractions of
these paths/relative positions and extended oscillation variables; equivalently keep all
the shifts as separate entries and use the deterministic dilation identifications. Window
extrema can use boundaries at deterministic nonjump times. These arguments also let us
keep tests at finitely then countably many bases, windows and cuts; joint independence of
tangent paths at different bases is not asserted.

\end{proof}

\smallskip
\noindent\textbf{Protected apertures and near-links.}\quad

Suppose on the finite event there is exactly one positive half, for example the \(h\)
coordinate toward the past. Look at the level \(H=h_t\). At \(t-r_y\) (round throughout)
it is strictly in the list, whose other end in every compact window is very far below
\(H\). Trace the strict inclusion of \(H\) backwards from there, to the initial wall if
uninterrupted, or to the last crossing from outside (including equality as outside). The
crossing step introduces it from a positive aperture. The orientation can be at either end
of the main list. The reverse procedure if the positive half is toward the future gives a
removal step, or the final survivor. A bounded loss/end convention at a crossing is
irrelevant unless the level is near an endpoint. Use reflected labels for the other
coordinate. This choice of event/disk/wall is the \emph{same} at all fixed nearby scale
shifts, asymptotically on the positive-half event (strict interior uniformly over the
intervening times), and its time is infinitely far in viewed units. For every fixed large
\(K\) keep, divided by \(P_y\), the least endpoint distance to \(H\) at the crossing (both
ends of the list in both states), or initial/final state as appropriate, and throughout
the subsequent (respectively preceding) sustained interval up to time \(t-Kr_y\)
(respectively back to \(t+Kr_y\)). Use extended limits at shifts and increasing \(K\),
e.g.~powers of \(R\) allowing vanishing rounding errors. The limit as \(K\to\infty\) is
either zero or infinite, by stationarity, stability of the chosen crossing, and
monotonicity. Bounded-factor time sandwiches suffice for this argument with rounded window
endpoints. In the infinite case there is positive clearance already at fixed positive
viewed distances as needed, by compact strict positivity. In the zero case we get a
near-link to a time \(q\) arbitrarily far in viewed units, at which one of \(h_q,b_q\)
(allow either state-copy) is within \(P_y\) of \(H\). Throughout the interval between
\(t\) and \(q\), both relevant endpoints stay on their respective required sides of \(H\)
(\(h\) above, \(b\) below) with an \(O(P_y)\) tolerance tight on this finite event. Indeed
there is strict persistence from the crossing to the unit time window, except possibly the
outer pre/post state used at that crossing where the near equality suffices, and core
oscillations in the unit window are finite. A last crossing at a degree-one or other
bounded loss without a true large removed aperture necessarily uses the near-link
alternative. These statements at large \emph{fixed} factors first, with high limiting
probability on the cases in question, can throughout be used at factors tending
sufficiently slowly to infinity.

\smallskip
\noindent\textbf{Ring transfer on a protected aperture.}\quad

Call the infinite-clearance alternative protected. A protected aperture containing \(H\)
has margins at its ends infinitely large at this scale. Through positive rings here the
relative transfers about \(H\), on bounded parts of this scale away from the aperture
endpoints, are linear unit-slope in the limit (with the specified cyclic orders); this
also holds at each fixed shifted scale. For detail all relevant past positives before the
endpoint strip have outer sizes \(\le C p\). Before tilt whole rings are independent
geometric runs given the outer sizes, with total tilt by \(f_l\); on comparable totals the
density factor is bounded, and the tilted atypical estimates for uniform-order deviations
on interval lengths of order at least the smallest tested positive scale beat any power of
\(p_*\). This follows by the exponential geometric sum bounds, even after the polynomial
reciprocal normalization, a union over slots and over main-path steps, then
\eqref{eq:original-19} or the event/sampling likelihood. In particular the aperture arcs
are very large on the fresh boundary also. At an initial wall we need no transfer to an
auxiliary disk.

\smallskip
\noindent\textbf{Accessibility and earlier helper interference.}\quad

For each protected \(J\), distances from the corresponding seed for \(H\) to the other
prescriptions there in cyclic boundary arclength, divided by \(P_y\), tend to zero or
infinity in these stationary extractions: their unscaled values are unchanged at fixed
scale shifts on the invariant event. Select the earliest seed in this zero-distance group
(thus offset \(o(P_y)\)). Before it, each earlier path either never interferes near this
arc or has an interference. More precisely take cyclic distances to the closest
consumption of an original slot by that path, again zero or infinity on the viewed scale.
If there is a zero earlier, charge to the \textbf{first} such path. When it started, an
arc about \(H\) with both margins infinitely large in these units was pending in its
polygon, and the path started far in boundary arclength from the center of that arc on
both sides. Otherwise such an arc is pending at the selected near helper itself (the
accessible case), and its start slot is available and near-centered. These implications
follow from untouched-arc persistence; in the interference case the arc must have been in
the charged starting polygon in order to be reached, and even if remote parts of the
circumference have been compressed, the long contiguous fresh arc separates its middle
from the earlier start. The same stationarity assertions hold if retaining these discrete
choices/tests at shifts jointly by subextraction; stability of the selected near helper
then holds on the protected event under those shifts.

An interference forces on at least one side of that earlier path the \emph{descending
ladder range} of its perimeter-coordinate increments from time zero to come within
\(O(P_y)\) of the corresponding depth \(x\) (distance down that side at start to the
middle label, measured in the transmitted arc in the starting polygon). Indeed at the
first consumption in a window of that radius the original slot is still persistent up to
its transfer to active or its deletion/pairing. To touch as active requires reaching its
label within the bounded endpoint convention; to consume as paired by a loss lands just
beyond it on the updating side. No earlier record on the reaching side can already be
deeper beyond the window, since that particular occurrence persisted. Thus a record depth
lies in the window with a bounded enlargement. Here \(x/P_y\to\infty\).

\smallskip
\noindent\textbf{The accessible helper tangent.}\quad

In the accessible case perform the same centered-path tangent tests on the chosen helper,
using forward time from 0 on its two sides. Its starting size is \(\gg P_y\). On the
finite-oscillation case it runs longer than every fixed viewed time by the balanced
stopping rule. Its alternatives are singular oscillation at the origin, one or more
positive halves, or two strict dips with continuity at the origin. The entropy proof for
increments is the same, conditioning first on the entire main word/seam geometry and
appropriate earlier helper word histories. The complete conditional entropy argument is
given in Lemma~\ref{end:record-null-transfer}. Given the main choices,
partition queries by their helper prescription index (there are only boundedly many
indices per disk, the partition need not inspect the interior). One can make the queried
path when possible or adjoin artificial iid increments, conditioning on the preceding
helper operations but not their unused fillings. Within each disk the windows for that
index use disjoint slots of the same augmented compact-likelihood walk, and between disks
the original new path inputs conditional on starting data are products. Initial unused
gaps also give ordinary position smoothing. The stability just proved and stationarity
give the core assertions as for the main path.

\begin{proposition}[Exclusion of uncharged passage configurations]\label{end:passage-exclusion}

The following possibilities have probability zero for typical-scale limits under the lift:

\begin{itemize}
\item
  all four main halves dip strictly (finite centered test);
\item
  just one positive, protected, and either initial wall, or accessible with the helper's
  two strict dips (finite centered tests).
\end{itemize}

\end{proposition}

The assertion concerns spatial passage and needs more than polarity of the
Cauchy coordinates. We prove it through buffered records, a comparison with
ordinary paired disks, and an incidence certificate that is valid in every
completion of those records.

\subsection{Buffered records and adaptive entropy}

If one of the possibilities in Proposition~\ref{end:passage-exclusion}
persists with positive probability, we can restrict to fixed windows and strict constant
buffers in it, with still positive limiting probability. Take the main outer cuts at times
\(-U_-,U_+\), and inner cuts at \(-\eta,\eta\), in viewed relative time, rounding to
steps. In the helper when used take cuts at \(\eta',U'\) measured from its start. The
following conditions can be imposed, reflecting either label/time description when
appropriate:

\begin{itemize}
\item
  In directions requiring dips, they occur strictly between inner and outer cuts and to
  depth more than \(10D\) below the corresponding centered value (initial value for
  helper). Here \(D>0\) can be small. Relative walk values in the windows are bounded by a
  fixed \(M\). Central-gap oscillations are much less than \(v\), where \(0<v\ll D\); we
  can ask, for example, for endpoint offsets less than \(v/1000\) and gap oscillations
  also that small, by taking the inner cuts small.
\item
  On a protected main direction, its outer cut is more than \(10D\) \textbf{above} center,
  and the main arc \(|\text{label}-H|\le 2D P_y\) is received unchanged from the
  designated aperture/initial wall up to that cut, or transmitted unchanged from that cut
  to the future loss aperture/final disk. It is strictly inside the appropriate range
  there, with room. None of the walk comparisons in and between retained main blocks
  requires mixing heights near the two ends across the length of the list (in the actual
  test they are separated by a diverging width in viewed units).
\item
  On the other side of a designated aperture the corresponding arc, in its actual slot
  order and with the ring transfer when needed, is intact until the prescribed helper
  segment, whose start has arclength offset much smaller than \(vP_y\) from the
  corresponding middle position. Its own list sides are separated toward the back far
  beyond the tested size bounds, and its run covers the helper blocks. We can use
  unit-slope relative accuracy as fine as needed on the actual ring interval. The future
  survivor case uses the same direct arc transfer type as a loss aperture. There is no
  helper at an initial wall.
\end{itemize}

Fix the finitely specified types of these conditions, including the coordinate choices.
The list conventions include orientation and boundary color. Use as local outputs:

\begin{enumerate}
\def\labelenumi{\arabic{enumi}.}
\item
  the main raw words on the two outer-to-inner blocks (outer on the first, inner on the
  second as initial times), with \emph{all} their detached raw contents and ring data, and
  the integer offsets of both coordinates at the two inner cuts relative to their values
  at \(t\);
\item
  at a positive designated \(J\), ring chunks on consecutive outer bases for each of the
  two label intervals extending from distance \(vP_y\) to \(D P_y\) on opposite sides of
  \(H\). They contain the corresponding inner runs with the exact local matching (use
  successive outer bases as binary-list cuts). For direct transfer or initial wall just
  use the corresponding interval of main labels. One can round inwards by bounded amounts;
\item
  the helper raw word and all its detached contents on the later block \([\eta',U']\) when
  used. Also use \textbf{two} integer offsets at its near cut: measure on each side
  relative to a fixed arc slot/vertex anchor in the corresponding transmitted interval
  (e.g.~the cut nearer \(vP_y\), using the corresponding inner-list corner at the ring).
  Continue its integer coordinate even past consumption, so this is initial oriented
  distance plus signed increments. Thus comparisons between helper operations and both arc
  pieces can use exact labels. They need not assume a known arclength shear between the
  two ring chunks. Which increasing arc direction is used on each helper side follows from
  the cyclic order and near-centered start.
\end{enumerate}

The central and exterior buffers above on the \emph{omitted} data will be side conditions
for transfer; keep strict slack, e.g.~allow oscillations a little larger but less than
\(v/5\) in the central gaps, origins displaced by \(<v/10\) on the helper seam, and still
positive excess in the exterior and width clearances. We can restrict observable
words/coordinates by the indicated bound, dips, protected-side value at its outer cut,
near-end offset bounds (on helper offsets use a small fixed tolerance, say \(v/200\),
about the positions predicted by unit-slope transfer and a centered start), and by
similarly accurate linear relative transfer in the retained ring chunks. Use strict slack
in these restrictions also. All omitted main starting/ending times for our bad-time test
lie outside the outer cuts and its grid-radius neighborhood is inside the inner cuts.

Denote this retained record by $\mathcal R$. It consists of the main
block words and their raw contents, the main inner-cut offsets, the two
seam chunks or direct-transfer label intervals, the helper block and its
two anchor offsets when present, and the finite local transfer and orientation tags. The local transfer
type is a positive ring, a direct ordinary-slot transfer, or an initial
wall, together with the protected-coordinate tag and orientation. The helper-selection
history may distinguish a loss aperture from a final survivor; their local
record type is the same direct transfer.
For a fixed record $\mathcal R$, let $\mathfrak C(\mathcal R)$ be the set
of its finite balanced-disk completions satisfying the omitted-data buffers
above. We seek a finite primal separator encoded by $\mathcal R$ that is
a good-time witness in every member of $\mathfrak C(\mathcal R)$.

The next two lemmas compare the law of this same record in two directions.
The entropy estimate transfers events of vanishing elementary-reference
probability to the lifted law. The synthetic experiment supplies the reverse
domination needed to show that failure of the universal separator property
has vanishing elementary-reference probability.

\begin{lemma}[Null transfer for complete buffered records]\label{end:record-null-transfer}

For these outputs the lift at a typical log scale transfers null probabilities from the
following elementary lattice references, with ordinary compact cutoffs or dummy symbols as
needed: raw iid words on the indicated blocks with their conditional detached contents;
independent ordinary geometric runs on the two designated ring intervals when present;
independent lattice offsets on bounded boxes with masses of order \(P_y^{-1}\) per scalar.
Here we only test the prescribed open-buffer regimes; there is no unconditioned domination
claimed for a microscopic tangent centered on a specially selected slot with no noisy
offsets.

\end{lemma}

\begin{proof}

\smallskip
\noindent\textbf{Conflict-free main and seam observations.}\quad

Here are details to include the raw contents in that entropy argument. Treat one base,
coordinate/crossing type and fixed window choices at a time. First handle the main
increments and main offsets by the separated annular sublattices already described. The
subsequent estimates on a given sublattice can condition on the \emph{full primitive main
walk} (and choices of bases) to decide which designated aperture to try by tracing the
level, without conditioning yet on its contents or rings. Query on the required
large-aperture ranges if outside the retained main time blocks of this test scale. Drop
other queries. Resolve overlap conflicts before further comparison: an aperture step in
one scale's query can occur in a queried main block for a different scale in the spaced
list, but has at most one such block index. Thus color the indices with a bounded number
of colors. Indeed the directed conflict graph has at most one outgoing edge per
index, so each connected component of its underlying simple graph has at most
one cycle and admits a coloring with three colors. Within one color there is no
such conflict. Initial/final boundary disks are treated without a step
conflict. Terminal survival queries refer to the actual final polygon, separate from all
the retained detached pieces. Repeated queries of the same positive seam here use disjoint
chunks around the same \(H\) by sufficient sublattice spacing.

\smallskip
\noindent\textbf{Conditional product laws and adaptive helper observations.}\quad

Conditional on the main word, detached pieces in retained main blocks have their raw
product law, and ring decoration at the designated positives, independently, can use
ordinary geometric sequences with the \(f_l\) tilt. This tilt has bounded log-density
above per queried ring on comparable totals. Noncomparable totals for these queries cost
negligible even in the positive part of the log-density expectation under
\eqref{eq:original-19}, since queried outer sizes are at least a constant times the
smallest scale and at most \(O(p)\); the untilted-law atypical estimates are exponential
in outer size up to polynomial factors (under the original \(f_l\) law as well). Disjoint
chunks at repeated scales thus cause only \(O(1)\) cost per index on average. Now, given
the whole applicable ring data, one can handle helper outputs separately for each possible
prescription index within disk. For that index, the prior sigma-field contains
the primitive main word, applicable ring and seam data, and earlier helper
word histories. It excludes unused map interiors and raw filling outputs of
other prescription indices. Condition on these preceding operations for those disks.
The starts, pending-arc sufficiency and transmissions of anchors can
then be chosen before those new paths; invalid starts/anchor queries can have dummy
outputs. For the same prescription in a disk the later word blocks queried are disjoint
across the spaced scales and have fresh inner intervening gaps. Use bounded path-law
likelihood once per such trajectory (iid extension and artificial fresh step contents
after stop if necessary), conditioning detached contents of its queried blocks by their
raw kernel. Thus both requested initial translations, even with shifts supplied by rings
or preceding path histories, use the joint gap local-mass bound \(C/P_y^2\). They need not
be centered by an exactly resampled source. Fillings inside a helper word block are not
further queried for this same prescription in that disk.

\smallskip
\noindent\textbf{Joint records and the entropy chain rule.}\quad

We combine the comparisons in their order of revelation. The needed entropy
inequality allows a prior to contain more information than the retained
record. Let $R$ be a retained record, let $Z$ be a prior with $R$ measurable
from $Z$, and let $U$ be newly queried outputs. Write $Q$ for the lifted law
and $P$ for the original disk law. If $\nu$ is the reference law of $R$ and
$K_R$ is a reference kernel for $U$ depending only on $R$, convexity and the
chain rule give
\begin{align*}
 H\bigl(Q_{R,U}\mid\nu(dR)K_R(dU)\bigr)
 &\le H(Q_R\mid\nu)+\E_Q H\bigl(Q(dU\mid Z)\mid K_R\bigr)\\
 &=H(Q_R\mid\nu)
  +\E_Q H\bigl(Q(dU\mid Z)\mid P(dU\mid Z)\bigr)
  +\E_Q\log\frac{dP(dU\mid Z)}{dK_R(dU)}.
\end{align*}
Thus one can forget the extra prior information after comparing the new
outputs jointly. The conditional entropy term is at most the full lift
entropy budget, by the chain rule. A new finite tag $T$ selected from $Z$
can be included by giving it an independent uniform reference; this adds
at most $\log |\mathcal T|$ per queried index. Here $\mathcal T$ denotes
its finite alphabet.

Apply this inequality first with $R$ equal to the retained primitive main
words and noisy offsets on one spaced scale list. Their product reference
and $O(N)$ entropy bound were established before revealing the full main
word. At this stage $Z$ contains that full primitive word and the base
choices. The designated aperture and the conflict color are determined
from $Z$. The new outputs are the main-block detached contents and the
queried ring chunks. Within one color, the designated aperture lies outside
all its retained main blocks; repeated seam queries use disjoint chunks.
The reference kernel therefore uses the original filling kernels given the
retained words and ordinary geometric runs on the queried chunks. It depends
only on the retained record and its finite presence/type tags. The $f_l$
ring tilts and their atypical-tail estimates give expected logarithmic cost
$O(N)$ for this entire conditional list.

For helpers, fix one prescription index at a time. There are only boundedly
many possible indices, and a complete local record uses just one of them.
The common retained record $R_0$ consists of the main and ring outputs already
compared. For index $j$, take a prior $Z_j$ containing the full primitive main
word, the applicable ring and seam data, the retained main-block contents,
and earlier helper \emph{word histories}. It excludes unused interiors and
raw filling outputs of other helper prescriptions. In particular, we do not
put earlier complete helper records into $R_0$: such a raw filling could
contain a polygon that the current helper will explore.

The new output $U_j$ is the record for prescription $j$, including its own
detached contents. Starts and transmission anchors are chosen from $Z_j$;
at each scale include the finite indicator $I_j$ that this prescription
supplies the accessible helper, and use a dummy output when $I_j=0$ or
the query is invalid. Accessibility and preceding-path interference or
availability choices made from this prior enter as finite tags here. Give these tags the independent
finite reference described above. The conditional comparison is made separately
for each $j$ and then summed over this bounded set of alternatives.

For a fixed prescription its queried blocks at the spaced scales are
disjoint and have fresh intervening gaps. The joint gap bound gives mass at
most $C/P_y^2$ for the two offsets, uniformly in the absolute anchor shifts
supplied by $Z_j$. Its reference kernel $K_j$ therefore uses independent
offsets on the prescribed bounded boxes, raw helper words, and their original
conditional filling kernels, with parameters determined only by the retained
local anchors and finite type. The sequential bounded path likelihood and
gap bounds give expected logarithmic cost $O(N)$ against the product of these
new record kernels. Applying the displayed inequality to $(R_0,Z_j,U_j)$
gives a joint entropy bound $O(N)$ for the main/ring record together with
this helper. It asserts no product law between records of different helper
prescriptions on the same disk.
These elementary references factor over scale indices, retaining the
within-record dependence of detached contents on their own raw words.
Superadditivity of entropy over the factors bounds the sum of the
one-record entropies for each masked record $(R_0,I_j,U_j)$ by $O(N)$.
Divide by the number of sampled scales, comparable to $N$. At one sampled
scale the actual complete record uses at most one helper prescription.
An elementary-reference null event for that record therefore has lifted
probability bounded by the sum of the corresponding masked-record
probabilities, over the boundedly many colors and indices $j$. Each tends
to zero by the entropy bound. This proves the claimed joint null transfer
without a joint law for all alternative helper payloads. No positive lower
bound on the size of an individual selected list is needed. In
particular, main-offset smoothing was used only in the initial term
$H(Q_R\mid\nu)$, never conditional on the full main word.


The probabilistic implication can be written explicitly. If \(Q\) is a
lifted record law, \(P\) its elementary reference, and \(H(Q\mid P)\le C\),
then for every event \(A\) with \(0<P(A)<1\), the binary entropy bound gives
\[
 Q(A)\log\frac1{P(A)}\le C+\log2.
\]
The averaged conditional entropy bound gives the same conclusion for the
uniformly sampled scale. Thus reference probabilities tending to zero force
lift probabilities to tend to zero. This also justifies the helper tangent
path absolute-continuity uses without raw contents.\qedhere

\end{proof}

\subsection{Ordinary synthetic experiments and universal incidence certificates}

\begin{lemma}[Reverse comparison on buffered records]\label{end:synthetic-domination}
On each of the strict regular subsets of records just described, there is an
ordinary paired-disks experiment, in perimeter units $P_y$, whose completions
satisfy the omitted-data buffers and whose projection onto the retained records
dominates a fixed positive multiple of the elementary lattice reference law.
The statement allows tighter inset offset and ring-accuracy bounds. Its constants
depend on the fixed buffers, not on the diverging units.
\end{lemma}

\begin{proof}

\smallskip
\noindent\textbf{Three preparations of the distinguished aperture.}\quad

We explain the reverse comparison on the strict regular subset of records just specified.
Use a synthetic \textbf{ordinary} paired-disks experiment in perimeter units \(P_y\), with
the main side of the required color a first central-side caller, its length in a
sufficiently large fixed compact positive range (smooth scale mixing), with ordinary fresh
source and target. Explore its targeted branch. Put a deterministic central time in it
with both blocks as above in exact step lengths, leaving preparation time before them;
request positive large side lengths throughout all prescribed main operations. Below, a
``large'' size in this experiment can be fixed as a sufficiently large constant in viewed
units. Restrict the initial split and size to appropriate large intervals with room.
Simulate the relative \(h,b\) directions with the corresponding two main list coordinates;
the absolute separation of those directions need not be the original one. Orders on slit
arcs and the two color types agree under the following choices:

\begin{itemize}
\item
  If no aperture is needed, prepare with small relative displacements before the first
  window. For the initial-wall case also do this: the protected past outer-cut height
  exceeds the center height by the stipulated margin but a bounded amount, so \(H\) in
  this realization is strictly on the ordinary original arc and protected as required.
\item
  For a past positive aperture, make a large positive jump on the synthetic
  \textbf{protected coordinate} in preparation, with magnitude much larger than the offset
  bound and helper window bounds, and keep other preparatory fluctuations small compared
  with the margin. Then the band about its \(H\) has arrived from that jump unchanged. The
  increasing pending order and reverse ring gluing there depend only on the aperture type,
  even if in the actual test the giant step reached \(H\) through the other end of the
  lifted list.
\item
  For a future loss aperture or final survivor, use a \textbf{negative} separation after
  the protected outer post-cut, on that same synthetic coordinate, dropping by a large
  amount across \(H\) with both margins. Keep intermediate fluctuations small. Arrange by
  a still larger preparatory positive gain \emph{before} the windows that this whole
  separating interval is strictly above the initial side height. It then carries no
  inherited caller targets, and the negative disk is reached for an actual fresh sweep.
  Indeed original target slots transmitted along this path remain on their original
  prefixes until left, and on the limit the corresponding loss circle has no contact with
  the original boundary by the strict stack rule (opposite side has no record linkage at
  that jump, also immediate by its Cauchy oscillations on both time sides simultaneously
  at the countably many jump times of the first coordinate under compact base-law
  domination). This is a compact canonical negative restart in the primary comparison.
  Main target positivity and sufficient size room on both lists can all be imposed. The
  local data at this loss use exactly the direct transfer type, regardless also of whether
  the actual path stopped as a survivor instead.
\end{itemize}

These preparations reproduce the retained record type, rather than the
entire balanced exploration. A positive seam is read through its ordered
outer bases and inner runs; that local order does not record which remote
end of the main list introduced it. A loss aperture and a final survivor
both transmit an unchanged interval of ordinary slots with its inherited
order and color. Their retained direct-transfer data are therefore identical.
Only those interval data and the stated buffers enter
$\mathfrak C(\mathcal R)$, which is why the synthetic negative restart may
represent either case.

\smallskip
\noindent\textbf{Ordinary sources, targets and exact noisy offsets.}\quad

When used, \(J\) now has its ordinary fresh source, which we restrict to a small
macroscopic tolerance interval about the desired label in boundary order, and an ordinary
target with enough room for its one helper block and terminal clearance outside the seam
band. Use the targeted branch from that source. No earlier helper operations have to be
simulated: the wide pending unaltered arc in the test is sufficient. Thus all synthetic
spatial tests are on ordinary compared primary data (typical additional target choices are
allowed).

The mass on synthetic completions satisfying the required omitted-piece buffers, projected
to these records, dominates a fixed positive multiple of the elementary references
restricted by the indicated observable bounds/dips, with tighter inset offset and accuracy
bounds if needed. To verify this, all macroscopic preparatory and separating moves above
can use open jump time/size intervals of positive raw probability with intervening small
oscillations, by iid Cauchy convergence or truncated-jump sampling. Word likelihoods on
both compared branches have the positive bounded ratios \eqref{eq:original-12} on the
prescribed ranges. Given sizes, seam chunks use independent geometric runs up to the
smooth total tilt; complete the omitted runs with the indicated local linear accuracy with
probability bounded below, including between seam anchors needed for positioning the
helper start. Their chunk positions on outer bases are determined by main labels, not by
conditioning on inner counts. Both main offsets in each \emph{time direction} are made by
exact lattice displacements in the two omitted small-gap halves. The helper offsets use
two coordinate displacements in its initial omitted gap after selecting a source slot from
an ordinary tolerance interval as above: in particular, any bounded shear through the
seam's omitted middle within the chosen small tolerance is absorbed. The restricted source
probability is bounded below given those main/ring data; start/root occurrences can use
their true bounded lattice conventions here. For offsets as inset above, net displacements
required in each such gap are small compared with \(vP_y\); local masses of order at least
\(P_y^{-2}\) per gap are available while keeping oscillations \(<vP_y/5\), say. 

One direct
constrained local bound is to reserve, in a gap of fixed positive relative length, one
positive move in a small size band and one correcting negative of comparable size on each
coordinate in separate time portions. The bands can be strictly inside the
small-oscillation budget but larger than the required tiny net-offset and source-shift
bounds (decrease those bounds beforehand). Suppress other moves of these sizes or larger
and bound all remaining sums/oscillations by a sufficiently small fixed tolerance, an
event of probability bounded below: truncating at a still smaller fixed relative size
gives this by the truncated second moment, with negligible truncated drift, and the finite
larger-jump exclusions cost a fixed positive factor. The correcting moves can now set both
offsets exactly by summing over their possible locations, using \eqref{eq:original-12} at
each required integer in a comparable-size range; specifying each value costs order
\(P_y^{-1}\) after its time sum. Distinguished moves don't have an unbounded overcount
with these exclusions. This proves the lower bound with stricter numerical margins as
necessary. 

In the main path, central values and thus seam positions can be computed from
the earlier prepared heights, first block word and the sampled relative offsets; the
past/future protective inequalities in the recipes then hold for any such recorded words.
Detached raw contents inside all word blocks have exactly the same conditional laws. This
proves the comparison without any lattice conditioning at zero tolerance on relative
pure-field or attachment-source positions.

\end{proof}

It remains to prove a deterministic property of successful synthetic
records: their retained incidences carry one separator valid throughout
$\mathfrak C(\mathcal R)$. We first obtain a separator from the ordinary
spatial comparison. We then show that its complete vertex stars and side
crossings are determined by the retained record. This second step is what
allows the same finite path to be used in every buffered completion.

\begin{lemma}[A good-time witness in every completion]\label{end:universal-certificate}
Along good subextractions of the ordinary synthetic experiment, each buffered
retained record $\mathcal R$ admits a simple primal circuit or a simple primal crosspath
capped through the external root face with cost $O(\mathcal B(P_y))$. Its exact
incidences and separation properties hold in every completion in
$\mathfrak C(\mathcal R)$. Consequently it is a
witness that the tested center is a good grid time. The assertion includes
initial walls with pinched boundary walks and transmission to a survivor.
\end{lemma}

\begin{proof}

\smallskip
\noindent\textbf{Closed near and far sets.}\quad

Here is the certificate of success on this synthetic law portion. Along good
ordinary-experiment subextractions use the homeomorphic spatial comparisons, the exact
quotient contact descriptions and the open-port upper comparisons in \(\mathcal B(P_y)\)
units. The main disk is denoted \(D_0\). Work on limits where the buffer conditions above
still hold with slack. In particular the main protected synthetic band is on a single
targeted side with absolute heights well inside the relevant positive regimes; crossings
and cuts here are compact observations before terminal. Define closed \textbf{near} sets
by including:

\begin{itemize}
\item
  the whole main gap, trunk and jumped contents born there, and the frontier labels at its
  two cuts on each side within \(2v\) in scaled units of that side's central height;
\item
  in the seam or initial-wall case the boundary band \(|\text{label}-H|\le v\) (and the
  paired loci on a ring, which are identified in the limit);
\item
  in \(J\) when used, the initial helper gap and its born contents, and frontier arcs at
  its ends within \(2v\) of the corresponding initial tip heights.
\end{itemize}
Small roundings have no effect. The \textbf{far} sets include the main time portions
outside the outer cuts with all their jumped contents except distinguished \(J\)
(do include its main pinch); on \(J\)'s boundary include the complement of the
open band of radius \(D\) about \(H\). Include also the helper portion after its
post-cut through terminal with all its born contents. Include the exterior of
\(D_0\) with its rim, \emph{except} in the initial-wall test where we work only
in \(\overline D_0\), including the rim portion outside the open band about
\(H\) as far. All these definitions use closures, equivalently the unions of
the respective closed trunk pieces, closed disks and arcs by the vanishing
diameters at accumulation. Unknown root-face exterior information in an
initial-wall test will not be needed.

\smallskip
\noindent\textbf{Exclusion of every near--far contact.}\quad

Near and far are disjoint. Within one chart every linkage out of a central gap label or
the indicated near frontier arcs would have to sustain a level within that side's narrow
indicated heights. A required dip bars this before the corresponding outer cut. On the
remaining protected main direction, backward (or forward) linkage instead before that cut
into far times is impossible by the strict transmission from (or to) the open seam band:
no near-level operations occur in the intervening far period, and levels on the opposite
side of the exterior crossing are strictly beyond the band at that crossing. In the
initial-wall case the linkage is to the open ordinary original arc with the same strict
protection. These uses include aperture endpoints/trunk tips: in the exact-contact
description folds cannot propagate a nontrivial linkage into another direction/level
across a real-time tip. For a frontier point use its pending boundary label by the same
rule. Here heights of near frontier labels and of all gap apertures, on any working side
required for such ties, differ by much less than the dip/protection margins from center.
An aperture boundary can only use its actual side linkages including the automatic
endpoint fold at its birth; a class through a real time has at most one nontrivial side
linkage among its copies. Thus, for example, following a pending near height to a tip
outside the gap cannot then change to a different height to continue past the
corresponding barrier. The outer/target conventions give no further such contacts by the
strict buffers and original target room. Far parts within either chart cannot contact the
opposite chart's near parts: such a contact on the distinguished circle could use only
near-band levels there, all protected on the main side as just described (and barred by
the opposite-time dip), and barred from helper post times by the two dips since the seam
order is matched, the source has small offset there and the near helper arcs/gap meet that
seam only well inside the band. Interior disks in each trunk description have no other
contacts. This proves disjointness, also from the stated outer rim sets. The far set is
connected: main earlier and later parts join to the rim by source and target, with the
source and target on the connected included rim portion; the included arc on \(J\) joins
at its pinch, and the helper post part at the remote target. The near set is connected as well. Its frontier arcs contain the
corresponding cut tips by the small gap-oscillation bounds, and each
jumped closure attaches to the central trunk. In the seam case the
protected center level links this main piece to the seam band; the
helper source, when present, lies in that band. These are the same exact
contact and persistence rules used above. Thus the whole near set lies
in one component of the complement of far.

\smallskip
\noindent\textbf{Separating corridors, circuits and crosspaths.}\quad

In the sphere case that component is simply connected with the near set in it compactly
inset. There is thus a compact Jordan annular buffer there away from both sets, with a
simple circle separating the central tip from far (e.g.~use Riemann disk exhaustion
enclosing all near points in the component). For the initial-wall case flatten \(D_0\) to
a closed round disk and double by reflection of the sets through its simple boundary.
Doubled far is still connected, doubled far and near disjoint. A complement component as
above either is in the open disk or touches the boundary and is symmetric. In the latter
case use disk exhaustion invariant under reflection (Riemann map at a fixed point of the
reflection); the circle runs transversely across the flattened rim at two points. Use just
its simple crossarc on our main side, extending at both endpoint ports transversely into
the exterior in the \emph{ordinary paired sphere}. Endpoints have margin on the
transferable part of the initial arc. All corridors can be taken arbitrarily thin with
buffers; choices/perturbations can avoid the finitely excluded points for port comparison.

Ordinary metric length topology and open-port upper comparison now give respectively
surrounding simple primal circuits (force stitching by transverse crossings in the annular
corridor), or primal paths in the crossarc corridor passing through \(D_0\) from exterior
port to exterior port, at costs bounded in \(\mathcal B(P_y)\) units along further
pathwise extractions. Indeed follow finite reference paths of finite length for these open
instructions approximating the requisite Euclidean passages, with endpoint and confinement
clearance. Only fixed limiting tests or their countable strict approximations are used. In
the latter case crop to a real subarc through the main lattice disk with endpoints on the
two separated rim-port lifts. Such a crossing portion exists since central compact parts
of the crossarc are in the main open disk, transitions to the complement require a rim
touch, and the paired/exterior/ordinary placements have the all-lift upper comparison.
Erase loops to make the arc simple; it stays in a narrow crossarc neighborhood connecting
the two ends. Real arcs can terminate in the appropriate half-edges at boundary slots.
Along the blue primal graph all internal intersections use its actual planar drawing.

\smallskip
\noindent\textbf{Exact pairings on all retained stars.}\quad

The required control of a vertex star is uniform in its degree. Indeed,
Proposition~\ref{cut:disk-spatial} supplies comparison maps whose maximal
flag-image diameter $\delta_n$ tends to zero. Every flag incident to a vertex
$v$ contains $v$, so the image of its entire closed star has diameter at
most $2\delta_n$. Finitely many star enlargements therefore remain in any
fixed open buffer eventually, even when the number of incident flags grows.

The finite certificate consists of the separator's ordered edge
occurrences, its endpoint sectors in the crosspath case, the full rotations
of its vertices and adjoining stars, and the retained tip movements through
this neighborhood. We must show that all these data can be read from
$\mathcal R$, independently of the member of $\mathfrak C(\mathcal R)$.

We spell out why these real boundaries are transferable at exact incidence and preserve
separation in \textbf{every} disk test with these same records and the omitted-data
buffers, for sufficiently advanced successful synthetic comparisons. At a compactly
buffered point of the circle/crossarc (on the main/queried side in a crossarc test), all
the stars of lifted sites and boundedly adjacent incidence data stay off near and far by
the all-flag upper comparisons, including the aperture, cut and frontage label
comparisons. The triangle pieces there must all be from retained blocks with their raw
attached disks/rings, or retained seam chunks. Iterating their pending-slot operation
records confirms all actual pairings there:

\begin{itemize}
\item
  Within one raw block use its own stack operation order and comparisons. Across the main
  central gap, slots which pass and are needed on both sides necessarily stay outside the
  unknown thresholds (otherwise their synthetic points are in the near pieces or the near
  frontier arcs), and use the same offsets on their side. A pairing created by an omitted
  operation itself would lie in its omitted content/tip. Active-slot and bridge-pair
  assignments occurring locally on opposite sides at a retained operation use each
  separate end displacement; the absolute width across the far middle of the list isn't
  needed.
\item
  Comparisons with an exceptional seam through the far transmission require labels
  strictly inside the protected radius by far avoidance, and outside the omitted seam
  middle by near avoidance. They use the literal labels, contiguous ring chunks if
  positive, and on the helper side the two anchor offsets with an untouched sufficiently
  wide initial arc passed through any prior operations and the near cuts (source and
  omitted initial fluctuations have inner clearance). Thus both gluing ends of actual
  pairs at retained operations are reproduced even when comparing a loss aperture here to
  actual transmission into a survivor. The helper block on its own again uses just its raw
  word/fillings. A direct gluing through an omitted helper operation at these sites is
  excluded by avoidance.
\item
  At the initial rim include actual root slots and their successive boundary sectors on
  the indicated open arc pieces. Pairings on the \(T\)-side and root-sector rotation
  incidences are all known by those same label comparisons. There is no need to determine
  the true neighboring triangulation exterior to \(T\) on the sphere.
\end{itemize}

\smallskip
\Needspace{4\baselineskip}
\noindent\textbf{Composed occurrence pairings and closure of vertex stars.}\quad

More explicitly, initial/final cuts of a retained word block may be read with arbitrary
formal unfilled stacks. The only tests then needing translations outside the block to
confirm a pair are tracing \emph{ordinary} input/output slots by unchanged indices across
intervening omitted steps or a split/ring to their actual retained origin and consumption.
New/consumed actives at steps inside the block have the raw bounded neighbor conventions,
and do not silently persist across omitted peeling steps. Between main blocks a pending
relevant slot has index near one of their working ranges, which remain far apart in all
realizations used here. If its comparison could use a gap threshold within the allowed gap
fluctuations, its same synthetic pending height at a gap endpoint belongs to the enlarged
near arcs (or the slot reaches a gap operation). Between a seam and block there is exact
transmission as above, and across a helper start it is the assigned \emph{side} of the
source/range relative to its chunk that matters, fixed by the margin. A pair incident to
an omitted terminal polygon, raw ring chunk or word operation would have at least one flag
or peeled bridge/tip in the omitted sets. These checks exhaust the pending-slot gluings;
at a protected vertex the already confirmed pairs close its entire synthetic rotation
cycle (with root-face sectors in the boundary case). The raw flag/sector occurrences are
distinct in every realization, so additional unknown identifications cannot change such a
closed star or merge nonidentical ones.

To elaborate, along a gluing traced through
several pieces, root bridges with two boundary occurrences in a raw filling are read with
their full pairing from that filling. At such a trace all intermediary occurrences
representing the same final edge are themselves at the protected sites. A bridge deletion
by the axis word there must likewise be recorded (both its ends at bounded incidence of
its tip). Unchanged-slot transfers with comparisons on the main blocks need translations
only near the respective end working ranges determined by the bounded word windows; these
ranges are disjoint with margin, so a transmission through the gap cannot silently
exchange them. In comparing to \(J\), original slots used on the transmitting band up to
the helper start are still pending consecutively in its starting polygon regardless of
changes beyond; on each of the two protected band pieces the comparisons from that side's
near-cut anchor to actual helper operations now use literal increments, even across raw
bridge moves. Thus it is occurrence paths of gluing, not just abstract final vertex
equalities, which are confirmed.

Complete paired sides at a site trace back along these
paths to their same origins: constructions glue successive unfilled slot sides, with
boundary-only bridge occurrences passed between (an interior final edge has its two filled
sides at the ends of the composed pairing; at the original wall include its original
root-face sector if present). There is no additional operation later merging already
closed edges. In the root case every root-face corner needed there is itself flanked by
the indicated comparable original boundary slots by all-occurrence rim comparison. The
spatial buffers allow boundedly many star enlargements before this exact confirmation. In
particular simplicity along the transferred path, endpoints as occurrences, and the
rotations of all adjoining paths are exact. This argument does not require a simple actual
boundary at an initial wall.

\smallskip
\noindent\textbf{No omitted visits and preservation of sides.}\quad

No omitted main tip movement can then revisit the certified boundary in the test, since
making a main move there or boundedly adjoining would pair a peeled slot or introduce a
ring at that position by an operation not in the data; all such incidence operations at
the boundary and its adjoining stars have already been confirmed with their retained
locations/orders. Within retained main blocks, visits and passage sides read by the tip
path through that neighborhood agree. Thus on either side of the central gap the parity of
sides switched across the boundary is the same; one may use the oriented collar of the
real path and local rotations, counting changes from one open side to the other, to avoid
transverse-general-position assumptions. Gap endpoint/outer-cut ambiguities are away from
this collar by buffer separation. At terminal transitions of the actual balanced path
outside the cuts the old active at the shared endpoint is included in the same exclusion.
For a crossarc cap, attach an arc through the root face of the main disk (away from the
real tip walk except possibly at its bounding endpoints); its end-sector germs and all
corresponding sides can be matched by the known boundary-label rotations.

\smallskip
\noindent\textbf{Odd side counts, including a pinched root boundary.}\quad

For the crosspath, the winding argument uses the following precise
property: the comparison from the capped discrete sphere to the doubled
limit disk has degree one, and the chosen near and far inset image points
each have a unique preimage in the main disk. The collar and cap avoid
those image points. We construct this comparison before applying parity.

For detail the side counts on both retained main blocks are odd. For a full circuit this
is the surrounding/separation property on the synthetic spatial comparison. In the
crossarc case it can also be checked as follows, allowing pinched root boundaries. Give
the root face artificially back to the synthetic main disk, thickening its real closure by
a collar in that face to obtain a simple closed topological disk, then cap (remove a
smaller simple open disk from the parametrized root-face interior, using its prime-end
boundary walk at the real sites). There are continuous degree-one comparisons from this
sphere into the doubled \(D_0\) on the good subextractions: on the real part take the
approximate spatial homeomorphisms retracted to \(\overline D_0\) from a small outer
Jordan-coordinate neighborhood, identity retraction inside; over the collar interpolate to
the ordered boundary parametrization using the uniform rim-label comparison (all widths
here can tend to zero in the images). On the capping disk use a homeomorphic extension to
the reflected disk. On a compact inset of the main open disk these comparisons have
ordinary homeomorphic degree-one lifts, since the complementary parts are spatially
outside with margin there. Extend the cropped real path through the collar at its two
separated slot-sector ports, joining through the cap along the reflected crossarc. Its
Jordan loop then projects within the annular buffer and winds once between the indicated
sides (the real path lies in a narrow crossarc corridor joining the two ends; backtracking
does not change this homotopy). The synthetic central time and representative times in the
pre- and post-outer pieces give the separated points on the respective sides. We can take
them with limits \emph{interior} to the main disk: deterministic compact-surviving trunk
times almost surely have no outer-rim contact by Cauchy oscillation and the quotient rule.
Allow extra surviving pre/post buffer times in the preparation if needed, or use a dense
deterministic sequence. Winding parity between their lifts follows from the degree-one
local property there: inset image positions at those tests have unique preimages where the
old homeomorphisms apply. If the Jordan loop did not separate those lifts, the opposite
source disk side would contract the projected loop away from both positions, impossible.
Since no omitted tip visits the real boundary or the cap between tests, this gives odd
switches across each retained block just as for a full circuit. Actual capping in a test
realization can use any simple root-face arc with the matched germs.

\smallskip
\noindent\textbf{The transferred good-time witness.}\quad

It follows that in the actual \(T_p\) test one open side contains the whole grid-radius
neighborhood at the center, the full tip ends are outside, and its time component ends on
both sides within the outer cuts. The latter reach is smaller than the allowed reach, and
the costs here are within the allowed budget. For full circuits near winding curves, use a
simple separating subcircuit in the narrow buffer; at capped paths the cap contributes no
cost. The certified transfers work simultaneously for any completions with the indicated
margins, not by continuity of a discrete distance in abstract word displacements.

\end{proof}

\begin{proof}[Proof of Proposition~\ref{end:passage-exclusion}]

Consequently on the regular record subset the elementary-reference mass of possible
failure (existence of a buffered test completion without this good-time witness) tends to
zero. Indeed otherwise the synthetic lower comparison would put nonvanishing mass on this
event with synthetic buffers. On unconditioned good synthetic subextractions with those
buffered limits it is impossible by Lemma~\ref{end:universal-certificate}.
Finite strict-slack choices/exhaustion suffice, and possible completions or finite witness
failures at lattice indices are measurable (equivalently take the countable finite
combinatorial records/maps). This reasoning is uniform in the sense needed for a varying
log scale: a sequence of offending scales/addresses would have the same diverging units
and synthetic-law subextractions, with cut roundings and length-unit ratios as above.
By Lemma~\ref{end:record-null-transfer}, the same event has vanishing lift
probability. This contradicts the chosen bad bases on the stated
positive-probability buffered regimes. This proves the passage assertion.

\end{proof}
